\begin{proof}[Proof of \cref{obj:lem:occupancy-cancellation}]
By \cref{obj:ass:iid}, \(Q=(P^{\mathrm{obs}})^{\otimes n}\). All conditioning below uses this product law.

\begin{enumerate}
\item \textbf{Conditional cell contributions.}
Use the cells \(C_j\) of \cref{obj:synth_11}, and define
\[
p_j=P_X(C_j),\qquad
\mathsf P_j=\frac{1}{p_j}
 P^{\mathrm{obs}}\big|_{\{X\in C_j\}},
\qquad j=1,\ldots,k.
\]
Each cell has Lebesgue length \(\delta\), including the final cell with its right endpoint. Integrating the density bounds in \cref{obj:ass:density} gives
\[
\frac{\delta}{2}\le p_j\le\frac{3\delta}{2}.
\]
Moreover,
\[
0<\delta=\frac{2h}{k}\le h\le\frac14,
\qquad
0<p_j\le\frac38<1.
\]
Thus both cell membership and its complement have positive probability, and \(\mathsf P_j\) is a probability law.

Write the observed records as
\[
D=(O_i)_{i=1}^n,\qquad O_i=(X_i,A_i,Y_i),
\]
and define the pair kernels on two observed records by
\[
\begin{aligned}
K_N\bigl((X,A,Y),(X',A',Y')\bigr)
  &=(A-A')(Y-Y'),\\
K_D\bigl((X,A,Y),(X',A',Y')\bigr)
  &=(A-A')^2.
\end{aligned}
\]
Let \(S_{N,j}(D)\) and \(S_{D,j}(D)\) denote the respective cell summands in \cref{obj:def:private-ratio}. For every label set
\(\mathcal J\subseteq\{1,\ldots,n\}\), define
\[
\mathcal E_{j,\mathcal J}
=
\left\{D:\{i:X_i\in C_j\}=\mathcal J\right\},
\qquad
\mathsf Q_{j,\mathcal J}
=
\frac{Q|_{\mathcal E_{j,\mathcal J}}}
     {Q(\mathcal E_{j,\mathcal J})}.
\]
The product law gives
\[
Q(\mathcal E_{j,\mathcal J})
=p_j^{|\mathcal J|}(1-p_j)^{n-|\mathcal J|}>0.
\]
Indeed, restricting a product measure to this membership rectangle and dividing by its mass gives
\[
\mathsf Q_{j,\mathcal J}
=
\bigotimes_{i=1}^n
\begin{cases}
\mathsf P_j,&i\in\mathcal J,\\[2pt]
\displaystyle
\frac{P^{\mathrm{obs}}|_{\{X\notin C_j\}}}{1-p_j},
&i\notin\mathcal J.
\end{cases}
\]
Consequently, any two distinct coordinates with labels in \(\mathcal J\) have joint law
\(\mathsf P_j\otimes\mathsf P_j\).

The pair representations of the cell summands follow from
\[
\begin{aligned}
\sum_{\substack{i,i'\in\mathcal J\\i<i'}}
 (A_i-A_{i'})(Y_i-Y_{i'})
&=
|\mathcal J|\sum_{i\in\mathcal J}A_iY_i
-\left(\sum_{i\in\mathcal J}A_i\right)
 \left(\sum_{i\in\mathcal J}Y_i\right),\\
\sum_{\substack{i,i'\in\mathcal J\\i<i'}}
 (A_i-A_{i'})^2
&=
|\mathcal J|\sum_{i\in\mathcal J}A_i
-\left(\sum_{i\in\mathcal J}A_i\right)^2.
\end{aligned}
\]
Here the second identity uses \(A_i^2=A_i\). On
\(\mathcal E_{j,\mathcal J}\), therefore,
\[
S_{\bullet,j}(D)=
\begin{cases}
0,&|\mathcal J|<2,\\[3pt]
\displaystyle
\frac{|\mathcal J|}{\binom{|\mathcal J|}{2}}
\sum_{\substack{i,i'\in\mathcal J\\i<i'}}
K_\bullet(O_i,O_{i'}),&|\mathcal J|\ge2,
\end{cases}
\qquad \bullet\in\{N,D\}.
\]
Define the eligible cell count by
\[
s_j(D)=\sum_{i=1}^n\mathbf1\{X_i\in C_j\},
\qquad
s_{j,\ge2}(D)=
\begin{cases}
0,&s_j(D)<2,\\
s_j(D),&s_j(D)\ge2.
\end{cases}
\]
For \(|\mathcal J|\ge2\), the sum has
\(\binom{|\mathcal J|}{2}\) pairs, each with the corresponding moment
\(\nu_j\) or \(d_j\). For the empty and singleton label sets, both contributions and the eligible count are zero. Hence, in every case,
\[
\begin{aligned}
\int S_{N,j}(D)\,d\mathsf Q_{j,\mathcal J}(D)
&=
\left(\int s_{j,\ge2}(D)\,
 d\mathsf Q_{j,\mathcal J}(D)\right)\nu_j,\\
\int S_{D,j}(D)\,d\mathsf Q_{j,\mathcal J}(D)
&=
\left(\int s_{j,\ge2}(D)\,
 d\mathsf Q_{j,\mathcal J}(D)\right)d_j.
\end{aligned}
\]
The membership events form a finite measurable partition. Multiplying these identities by their event probabilities and summing gives
\[
\begin{aligned}
\int S_{N,j}(D)\,dQ(D)
&=\left(\int s_{j,\ge2}(D)\,dQ(D)\right)\nu_j,\\
\int S_{D,j}(D)\,dQ(D)
&=\left(\int s_{j,\ge2}(D)\,dQ(D)\right)d_j.
\end{aligned}
\]
The bounded kernels and the finite sample size ensure integrability throughout.
% lean: CausalSmith.Stat.PrivateCateRoughdesign.conditional_cell_contribution_expectation
% lean: CausalSmith.Stat.PrivateCateRoughdesign.cell_membership_setIntegral_normalize
% lean: CausalSmith.Stat.PrivateCateRoughdesign.integral_cell_membership_partition

\item \textbf{The occupancy weight.}
Pointwise, the eligible cell count counts precisely those cell members with a distinct partner:
\[
s_{j,\ge2}(D)
=
\sum_{i=1}^n
\mathbf1\left\{
X_i\in C_j,\ 
\exists i'\ne i:\ X_{i'}\in C_j
\right\}.
\]
For a fixed \(i\), the event that it belongs to the cell and has a partner is the cell-membership event minus the event that it is the sole member. Independence gives
\[
\begin{aligned}
Q\left\{X_i\in C_j,\ 
 \exists i'\ne i:\ X_{i'}\in C_j\right\}
&=
p_j-p_j(1-p_j)^{n-1}\\
&=p_j\left[1-(1-p_j)^{n-1}\right].
\end{aligned}
\]
Integrating the finite indicator sum therefore yields
\[
\int s_{j,\ge2}(D)\,dQ(D)
=
np_j\left[1-(1-p_j)^{n-1}\right]
=w_j.
\]
% lean: CausalSmith.Stat.PrivateCateRoughdesign.iid_cell_partner_count_expectation

\item \textbf{Population moment identities.}
By the definitions of the statistics,
\[
S_N(D)=\sum_{j=1}^k S_{N,j}(D),
\qquad
S_D(D)=\sum_{j=1}^k S_{D,j}(D).
\]
Linearity of integration, followed by the preceding two steps, gives
\[
\overline N
=\sum_{j=1}^k\int S_{N,j}(D)\,dQ(D)
=\sum_{j=1}^k w_j\nu_j,
\qquad
\overline D
=\sum_{j=1}^k\int S_{D,j}(D)\,dQ(D)
=\sum_{j=1}^k w_jd_j.
\]
% lean: CausalSmith.Stat.PrivateCateRoughdesign.population_cell_sum
% lean: CausalSmith.Stat.PrivateCateRoughdesign.occupancy_cancellation

\item \textbf{The conditional denominator floor.}
For each cell, define its conditional covariate law and mean treatment probability by
\[
\mathsf G_j=\frac{P_X|_{C_j}}{p_j},
\qquad
\overline e_j
=\int A\,d\mathsf P_j
=\frac{1}{p_j}\int_{C_j}e_P(x)\,dP_X(x)
=\int e_P(x)\,d\mathsf G_j(x).
\]
The middle equality is the defining conditional-probability identity for the selected propensity, restricted to the cell and normalized by \(p_j\). Averaging \cref{obj:ass:overlap} gives
\[
\frac14\le\overline e_j\le\frac34.
\]
For two independent binary treatments,
\[
(A-A')^2=A+A'-2AA'.
\]
Integration under \(\mathsf P_j\otimes\mathsf P_j\) thus gives
\[
d_j=2\overline e_j(1-\overline e_j).
\]
Finally,
\[
d_j-\frac38
=
2\left(\overline e_j-\frac14\right)
 \left(\frac34-\overline e_j\right)\ge0.
\]
Therefore \(d_j\ge3/8\).
% lean: CausalSmith.Stat.PrivateCateRoughdesign.conditional_cell_treatment_overlap
% lean: CausalSmith.Stat.PrivateCateRoughdesign.independent_binary_denominator_moment
% lean: CausalSmith.Stat.PrivateCateRoughdesign.conditional_cell_denominator_lower

\item \textbf{The conditional bias certificate.}
First express the observed pair moments through the selected regressions. Restriction to the cell and the defining arm-mean identities give
\[
\begin{aligned}
\int AY\,d\mathsf P_j
&=\int e_P(x)\mu_{1,P}(x)\,d\mathsf G_j(x),\\
\int (1-A)Y\,d\mathsf P_j
&=\int (1-e_P(x))\mu_{0,P}(x)\,d\mathsf G_j(x),\\
\int Y\,d\mathsf P_j
&=\int\left[(1-e_P(x))\mu_{0,P}(x)
             +e_P(x)\mu_{1,P}(x)\right]\,d\mathsf G_j(x).
\end{aligned}
\]
For example, the treated covariate submeasure has density \(e_P\) with respect to \(P_X\); integrating its selected outcome regression over \(C_j\) gives the first identity after division by \(p_j\). The control submeasure has density \(1-e_P\), giving the second, and their sum gives the third.

Expanding the numerator kernel and factoring the independent-record integrals yields
\[
\nu_j
=
2\left[
\int AY\,d\mathsf P_j
-\left(\int A\,d\mathsf P_j\right)
 \left(\int Y\,d\mathsf P_j\right)
\right].
\]
Define, for \(x,x'\in[0,1]\),
\[
\begin{aligned}
\mathcal K_{N,P}(x,x')
&=
e_P(x)(1-e_P(x'))
 \bigl(\mu_{1,P}(x)-\mu_{0,P}(x')\bigr)\\
&\quad+
e_P(x')(1-e_P(x))
 \bigl(\mu_{1,P}(x')-\mu_{0,P}(x)\bigr),\\
\mathcal K_{D,P}(x,x')
&=
e_P(x)(1-e_P(x'))
+e_P(x')(1-e_P(x)).
\end{aligned}
\]
Factoring the four terms in the integral of \(\mathcal K_{N,P}\) gives
\[
\begin{aligned}
&\int\mathcal K_{N,P}\,d(\mathsf G_j\otimes\mathsf G_j)\\
&\quad=
2\left[
\int e_P\mu_{1,P}\,d\mathsf G_j
-\overline e_j
 \left\{
 \int(1-e_P)\mu_{0,P}\,d\mathsf G_j
 +\int e_P\mu_{1,P}\,d\mathsf G_j
 \right\}
\right]
=\nu_j.
\end{aligned}
\]
Likewise, factoring the two denominator terms gives
\[
\int\mathcal K_{D,P}\,d(\mathsf G_j\otimes\mathsf G_j)
=2\overline e_j(1-\overline e_j)=d_j.
\]
% lean: CausalSmith.Stat.PrivateCateRoughdesign.conditional_cell_numerator_design_pair
% lean: CausalSmith.Stat.PrivateCateRoughdesign.conditional_cell_denominator_design_pair

For \(x,x'\in C_j\), the cell endpoints give
\[
|x-x_0|\le h,\qquad
|x'-x_0|\le h,\qquad
|x-x'|\le\delta.
\]
These inequalities include the final cell's right endpoint. The public radius is \(L=3\) by \cref{obj:synth_4}. Hence
\cref{obj:ass:contrast-lipschitz} gives
\[
|\tau_P(x)-\theta(P)|
\le3|x-x_0|\le3h,
\qquad
|\tau_P(x')-\theta(P)|\le3h.
\]
Similarly, \cref{obj:ass:propensity-holder,obj:ass:control-holder} give
\[
\begin{aligned}
&\left|
(e_P(x)-e_P(x'))
(\mu_{0,P}(x)-\mu_{0,P}(x'))
\right|\\
&\quad=
|e_P(x)-e_P(x')|\,
|\mu_{0,P}(x)-\mu_{0,P}(x')|\\
&\quad\le
(3\delta^{1/10})(3\delta^{1/10})
=9\delta^{1/5}.
\end{aligned}
\]
The overlap bounds also imply
\[
\begin{aligned}
\mathcal K_{D,P}(x,x')-\frac38
&=
\left(e_P(x)-\frac14\right)
\left(\frac34-e_P(x')\right)\\
&\quad+
\left(\frac34-e_P(x)\right)
\left(e_P(x')-\frac14\right)
\ge0.
\end{aligned}
\]
Both coefficients \(e_P(x)(1-e_P(x'))\) and
\(e_P(x')(1-e_P(x))\) are nonnegative.

Using \(\mu_{1,P}=\mu_{0,P}+\tau_P\), direct expansion gives
\[
\begin{aligned}
\mathcal K_{N,P}(x,x')
&=
(e_P(x)-e_P(x'))
(\mu_{0,P}(x)-\mu_{0,P}(x'))\\
&\quad+
e_P(x)(1-e_P(x'))\tau_P(x)
+e_P(x')(1-e_P(x))\tau_P(x').
\end{aligned}
\]
Subtracting \(\theta(P)\mathcal K_{D,P}(x,x')\), applying the triangle inequality, and using the preceding estimates yields
\[
\begin{aligned}
&|\mathcal K_{N,P}(x,x')
      -\theta(P)\mathcal K_{D,P}(x,x')|\\
&\quad\le
9\delta^{1/5}
+3h\,e_P(x)(1-e_P(x'))
+3h\,e_P(x')(1-e_P(x))\\
&\quad=
9\delta^{1/5}+3h\,\mathcal K_{D,P}(x,x')\\
&\quad\le
(24\delta^{1/5}+3h)\mathcal K_{D,P}(x,x')
=b\,\mathcal K_{D,P}(x,x').
\end{aligned}
\]
The last inequality uses
\[
9\delta^{1/5}
\le24\delta^{1/5}\mathcal K_{D,P}(x,x'),
\]
which follows from \(\mathcal K_{D,P}(x,x')\ge3/8\).

The probability law \(\mathsf G_j\otimes\mathsf G_j\) is supported on
\(C_j\times C_j\). Therefore,
\[
\begin{aligned}
|\nu_j-\theta(P)d_j|
&=
\left|\int
 \bigl(\mathcal K_{N,P}-\theta(P)\mathcal K_{D,P}\bigr)
 \,d(\mathsf G_j\otimes\mathsf G_j)\right|\\
&\le
\int
 \left|\mathcal K_{N,P}-\theta(P)\mathcal K_{D,P}\right|
 \,d(\mathsf G_j\otimes\mathsf G_j)\\
&\le
b\int\mathcal K_{D,P}\,d(\mathsf G_j\otimes\mathsf G_j)
=bd_j.
\end{aligned}
\]
% lean: CausalSmith.Stat.PrivateCateRoughdesign.cell_cross_treatment_certificate
% lean: CausalSmith.Stat.PrivateCateRoughdesign.conditional_design_pair_bias
% lean: CausalSmith.Stat.PrivateCateRoughdesign.conditional_cell_bias_bound

\item \textbf{Lower partner-probability bound.}
Since \(n\ge2\) and \(p_j\ge\delta/2\),
\[
n-1\ge\frac n2,
\qquad
(n-1)p_j\ge\frac{n\delta}{4}.
\]
The inequality \(1-p_j\le e^{-p_j}\), with \(1-p_j\ge0\), gives
\[
(1-p_j)^{n-1}
\le e^{-(n-1)p_j}
\le e^{-n\delta/4}.
\]
To compare the rational envelope with the truncated occupancy scale, split according to \(n\delta\). If \(n\delta\le1\), then
\[
\frac{\min\{1,n\delta\}}8(4+n\delta)
=\frac{n\delta(4+n\delta)}8
\le n\delta.
\]
If \(n\delta>1\), then
\[
\frac{\min\{1,n\delta\}}8(4+n\delta)
=\frac{4+n\delta}{8}
\le n\delta.
\]
Since \(4+n\delta>0\), both cases yield
\[
\frac{\min\{1,n\delta\}}8
\le\frac{n\delta}{4+n\delta}.
\]
Furthermore,
\[
e^{n\delta/4}\ge1+\frac{n\delta}{4}
\quad\Longrightarrow\quad
e^{-n\delta/4}\le
\frac{1}{1+n\delta/4}.
\]
Combining these estimates in order gives, for every cell,
\[
\begin{aligned}
\frac{\min\{1,n\delta\}}8
&\le\frac{n\delta}{4+n\delta}
=\frac{n\delta/4}{1+n\delta/4}\\
&\le1-e^{-n\delta/4}
\le1-(1-p_j)^{n-1}.
\end{aligned}
\]
% lean: CausalSmith.Stat.PrivateCateRoughdesign.occupancy_partner_lower

\item \textbf{Upper partner-probability bound.}
Bernoulli's inequality and nonnegativity give
\[
(1-p_j)^{n-1}\ge1-(n-1)p_j,
\qquad
(1-p_j)^{n-1}\ge0.
\]
Thus
\[
1-(1-p_j)^{n-1}
\le\min\{1,(n-1)p_j\},
\qquad
(n-1)p_j\le\frac32n\delta.
\]
If \(n\delta\le1\), these inequalities imply
\[
1-(1-p_j)^{n-1}
\le\frac32n\delta
=\frac32\min\{1,n\delta\}.
\]
If \(n\delta>1\), they imply
\[
1-(1-p_j)^{n-1}
\le1\le\frac32
=\frac32\min\{1,n\delta\}.
\]
Hence the same upper envelope applies to every cell.
% lean: CausalSmith.Stat.PrivateCateRoughdesign.occupancy_partner_upper

\item \textbf{Total occupancy bounds.}
Summing the cell-mass bounds and using \(k\delta=2h\) gives
\[
h=\frac{k\delta}{2}
\le\sum_{j=1}^k p_j
\le\frac{3k\delta}{2}=3h.
\]
Define the total occupancy weight by
\[
W(P)=\sum_{j=1}^k w_j.
\]
Multiplication of the lower partner envelope by \(np_j\ge0\), followed by summation, yields
\[
\begin{aligned}
W(P)
&\ge
\sum_{j=1}^k np_j\,
 \frac{\min\{1,n\delta\}}8\\
&=
\frac{n\min\{1,n\delta\}}8\sum_{j=1}^k p_j
\ge\frac{nh\min\{1,n\delta\}}8
=\frac{\mathcal A}{8}.
\end{aligned}
\]
The upper partner envelope similarly yields
\[
\begin{aligned}
W(P)
&\le
\sum_{j=1}^k np_j\,
 \frac32\min\{1,n\delta\}\\
&=
\frac32n\min\{1,n\delta\}\sum_{j=1}^k p_j
\le\frac92nh\min\{1,n\delta\}
=\frac{9\mathcal A}{2}.
\end{aligned}
\]
% lean: CausalSmith.Stat.PrivateCateRoughdesign.occupancy_comparisons

\item \textbf{Denominator floor and ratio bias.}
Since \(0\le1-p_j\le1\), every weight satisfies
\[
w_j=np_j\left[1-(1-p_j)^{n-1}\right]\ge0.
\]
The population denominator identity and the cellwise lower bound therefore imply
\[
\overline D
=\sum_{j=1}^k w_jd_j
\ge\frac38\sum_{j=1}^k w_j
=\frac38W(P)
\ge\frac{3\mathcal A}{64}
=d_0,
\]
where the formula for \(d_0\) is supplied by
\cref{obj:def:private-ratio}. Because \(n>0\), \(h>0\), and \(\delta>0\),
\[
\mathcal A=nh\min\{1,n\delta\}>0,
\qquad
d_0=\frac{3\mathcal A}{64}>0.
\]
In particular, \(\overline D>0\). The common nonnegative weights also preserve the bias certificates:
\[
\begin{aligned}
|\overline N-\theta(P)\overline D|
&=
\left|\sum_{j=1}^k
w_j\bigl(\nu_j-\theta(P)d_j\bigr)\right|\\
&\le
\sum_{j=1}^k w_j
 \left|\nu_j-\theta(P)d_j\right|\\
&\le
\sum_{j=1}^k w_jbd_j
=b\overline D.
\end{aligned}
\]
Finally, division by the positive denominator gives
\[
\left|\frac{\overline N}{\overline D}-\theta(P)\right|
=
\frac{|\overline N-\theta(P)\overline D|}{\overline D}
\le b.
\]
% lean: CausalSmith.Stat.PrivateCateRoughdesign.occupancy_ratio_assembly
\end{enumerate}
\end{proof}
